% !TEX root = ../../main.tex
% Author-requested recall. Keep the full prime-root proof in its shared solution.
\begin{pendingreview}
\section{Divisibility, primes, and greatest common divisors}
\label{sec:appendix:divisibility}

\subsection{Divisibility and multiples}
For integers $d\ne0$ and $a$, the notation $d\mid a$, read
``$d$ divides $a$,'' means
\[
  a=md\quad\text{for some }m\in\Z.
\]
When $d>0$ and $a\ge0$, this says that $m$ groups of size $d$ make up
$a$. For example, $3\mid12$ because $12=4\cdot3$.
The multiplier $m$ may also be zero or negative. In particular,
$3\mid0$ because $0=0\cdot3$, and $3\mid-12$ because
$-12=(-4)\cdot3$.

For $a\ne0$, we can rearrange the defining equation as $d=a/m$ with
a nonzero integer $m$. The product form $a=md$ also handles $a=0$,
when $m=0$. The notation $d\nmid a$ means that there is no integer
multiplier. Thus $3\nmid10$, since $10/3\notin\Z$.
\index{divisibility}\index{multiple}

\subsection{Greatest common divisors and lowest terms}
For integers $a,b$ that are not both zero, their \emph{greatest common
divisor}, written $\gcd(a,b)$, is the largest positive integer dividing
both of them. For example, the positive common divisors of $12$ and
$18$ are $1,2,3,6$, so
\[
  \gcd(12,18)=6.
\]
The integers $a,b$ are \emph{coprime}, or \emph{relatively prime},
when $\gcd(a,b)=1$. If a prime $p$ divides both $a$ and $b$, it is a
positive common divisor, so
\[
  \gcd(a,b)\ge p>1.
\]

Every rational number can be written as $a/b$ with $a,b\in\Z$,
$b>0$, and $\gcd(a,b)=1$. This is a representation in \emph{lowest
terms}, obtained by dividing numerator and denominator by their GCD.
Being rational, or even being a prime integer, does not make every
representation reduced. For example,
\[
  3=\frac31=\frac62,\qquad
  \gcd(3,1)=1,\quad\gcd(6,2)=2.
\]
We choose a lowest-terms representation when an argument needs
coprime numerator and denominator.
\index{greatest common divisor}\index{coprime integers}\index{lowest terms}

\subsection{Prime divisors and Euclid's lemma}
A \emph{prime} is an integer $p>1$ whose only positive divisors are
$1$ and $p$. Euclid's lemma, proved in
\autoref{thm:analysis-i:euclid-prime-divisors}, states that for
$a,b\in\Z$,
\[
  p\mid ab\quad\Longrightarrow\quad
  p\mid a\quad\text{or}\quad p\mid b.
\]
In particular, $p\mid a^2$ implies $p\mid a$, since $a^2=a\cdot a$.
Repeated application gives $p\mid a^k\Longrightarrow p\mid a$ for
every positive integer $k$.
The prime hypothesis is essential. For instance, $6\mid2\cdot3$,
but $6$ divides neither $2$ nor $3$.
\index{prime number}\index{Euclid's lemma}

\subsection{How the prime-root contradiction uses GCD}
In \autoref{ex:hw01:prime-roots}, $p$ is prime and we suppose for a
contradiction that $\sqrt p$ has a lowest-terms representation,
\[
  \sqrt p=\frac ab,\qquad
  a,b\in\Z,\quad b>0,\quad\gcd(a,b)=1.
\]
The \hyperref[sol:hw01:prime-roots]{full solution} uses Euclid's lemma
to establish both $p\mid a$ and $p\mid b$. These conclusions force
\[
  \gcd(a,b)\ge p>1,
\]
contradicting the choice $\gcd(a,b)=1$. Rationality would have allowed
a lowest-terms representation, so this contradiction proves that
$\sqrt p$ is irrational. Throughout the argument, $p$ is prime;
the rationality of $\sqrt p$ is the assumption being contradicted.
\end{pendingreview}
